\documentclass[11pt,a4paper]{amsart}
\usepackage[T1]{fontenc}
\usepackage[margin=1in]{geometry}
\usepackage{amssymb}
\usepackage[hidelinks]{hyperref}
\newtheorem{theorem}{Theorem}
\DeclareMathSizes{8}{8}{8}{8}
\DeclareMathSizes{9}{9}{8}{8}
\DeclareMathSizes{10.95}{10.95}{8}{8}
\title{Core verification with at most three mutual friends}
\date{}
\begin{document}
\begin{abstract}
Given a partition of agents whose mutual friendships form a graph of maximum degree three, deciding core stability and deciding strict core stability are both coNP-complete under friend-oriented preferences without neutrals. A single reduction from cubic Vertex Cover produces a strictly blocking coalition in every yes instance and excludes even weakly blocking coalitions in every no instance. Together with the polynomial algorithms at degree two, this settles the degree threshold asked about in Simola's thesis.
\end{abstract}
\maketitle
\section{Preferences and the threshold}
Let $G$ be a finite simple graph. Neighbours are friends and all other distinct agents are enemies. An agent $v$ compares coalitions $S\ni v$ by the lexicographic pair $(d_S(v),-|S|)$, where $d_S(v)$ counts friends in $S$. For a given partition write $o_v$ for the number of friends and $p_v$ for the size of its block. A blocking coalition strictly improves every member's outcome. A weakly blocking coalition weakly improves every member and strictly improves at least one. Core and strict-core verification ask whether the respective coalitions are absent.

\begin{theorem}
For each fixed maximum-degree bound $d\ge3$, both verification problems are coNP-complete for mutual friendships without neutrals. For $d\le2$ both are polynomial-time decidable.
\end{theorem}

Simola~\cite[Sections 5.2 and 5.5]{simola} gives the polynomial case, recalls hardness at four friends and asks about three. Chen et al.~\cite{chen} discuss the missing case using degree six in the symmetric setting. There is a discrepancy between the neighbourhood-union definition in their text and the arc counts used for their thresholds. In the full version, Appendix~A.4 reduces the tractable case to paths and cycles, while Appendix~A.2 uses vertices with four mutual friends in the hardness construction. We count distinct mutual friends throughout and use the thesis, Proposition~5.2(iii), for the degree-two statement. No planarity restriction is imposed here.

\section{Construction}
A coalition is a polynomial-size certificate of either failure. For hardness take a simple cubic graph $H$ with $n\ge4$ vertices, $m$ edges and cover budget $0\le k\le n$. Vertex Cover remains NP-hard on these inputs~\cite{ak}. Alternatively, starting with maximum degree three, fill each unit of missing degree at $v$ with a fresh $K_4$ whose edge $ab$ is subdivided by $z$, joining $vz$. If $z$ is absent a cover needs $a,b$ and a third vertex; if $z$ is present it still needs two other vertices. Selecting $z$ and the two original vertices different from $a,b$ always suffices. Thus each unit raises the minimum by exactly three.

Orient each input edge for naming. Put $W=2m+2$, $L=W+6$ and $F=kL+2m+1$. All friendships are exactly those below. First make one partition block $C$.
\begin{itemize}
\item Give $C$ a cycle with $m+2n$ vertices, one $b_e$ per edge and two $b_x,b'_x$ per vertex. Give it also a path of $F$ vertices, called the fallback. For each input edge $e$ add a terminal $t_e$ adjacent to $b_e$ and to its own distinct fallback vertex.
\item For each $x$ add a gate $a_x$ and helper $j_x$, with edges $b_xa_x,a_xj_x,j_xb'_x$. Write $B$ for the vertices of $C$ outside the fallback.
\end{itemize}
The remaining vertices have the following blocks.
\begin{itemize}
\item For each edge $e=(x,y)$ make an adjacent pair $u_e,v_e$ its own block and join $u_e$ to $t_e$.
\item For each $x$ take two disjoint triangle blocks joined by one bridge. Of the four vertices not incident to the bridge, three are distinct ports for the three edges at $x$. Join each port to $u_e$ if $x$ is the tail, or $v_e$ if $x$ is the head. Connect the fourth vertex to $a_x$ by a path with $W$ new internal vertices, partitioned into consecutive pairs. The two triangles and those $W$ vertices comprise a gadget of size $L$.
\end{itemize}
The construction has $O(n^2)$ vertices and maximum degree three. Every cycle vertex, fallback vertex and helper has all its friends in $C$; each gate or terminal has two friends in $C$. Outside $C$ each pair member starts with one friend and each triangle member with two.

\section{Which coalitions can block?}
Call $S$ acceptable if all its members weakly prefer it. An agent in $S$ needs at least $o_v+\mathbf1_{\{|S|>p_v\}}$ friends in $S$, and needs $o_v+\mathbf1_{\{|S|\ge p_v\}}$ to prefer it strictly. In particular a vertex with all its friends in its own block forces all those friends and caps the size of $S$ at $p_v$.

Meeting the cycle therefore forces its entirety and then all of $B$. Meeting the fallback forces the whole path and every terminal. Meeting a helper forces the cycle. A gate needs two friends, so forces its cycle vertex or its helper. A terminal needs two, so forces either its cycle vertex or the fallback. Consequently an acceptable coalition meeting $C$ has size at least $\min(m+2n,F)>3$.

Outside $C$ there is no weak block of size at most three. Singletons lose an old friend. In a pair both members would have to be old pair agents with one friend each, leaving both indifferent. An acceptable triple must be a triangle. In the outside-induced graph the edges joining different blocks form a matching, as each outside agent has at most one such neighbour there. A triangle using different blocks would give some vertex two such neighbours. Thus only old triangles qualify, and all their members are indifferent.

In a larger acceptable coalition every included triangle vertex needs all three neighbours, and each included tether vertex needs both path neighbours. Either kind forces its entire gadget, the gate and all three switches attached to its ports. An included $v_e$ needs $u_e$ and its head gadget. An included $u_e$ needs two friends, and at most one is its terminal; hence it forces either its tail gadget or $v_e$ and its head gadget. Thus any outside member forces a gate and the cycle. If there is no outside member, meeting the fallback forces a terminal and that terminal needs its cycle neighbour. In particular every weak block contains $B$, has size at most $|C|$, and avoids the fallback: inclusion of the path would give all of $C$, leaving only $S=C$, which improves nobody.

\section{Equivalence}
In any weak block the terminal of each edge has lost its fallback friend and therefore requires $u_e$. At least one endpoint gadget is included for each edge. These full gadgets are disjoint, so the set $X$ of included gadgets is a vertex cover satisfying
\[
 |X|L\le |S|-|B|\le F=kL+2m+1<(k+1)L .
\]
Conversely, take any vertex cover $X$ with $|X|\le k$. Form $S$ from $B$, the gadgets of $X$, every $u_e$, and $v_e$ precisely when its head belongs to $X$. Then $|S|\le |B|+kL+2m<|C|$. Cycle vertices and helpers retain all their friends; gates retain both old friends, and each terminal substitutes $u_e$ for its fallback friend. They all strictly prefer $S$. Every chosen triangle has all its required external neighbours, every tether vertex has both path neighbours, and every $u_e$ has its terminal and either its tail port or $v_e$ by the cover property. Each chosen $v_e$ has $u_e$ and its head port. Every outside agent therefore gains a friend. The coalition strictly blocks. This proves hardness for the complements of both verification problems and, with membership and the cited polynomial case, the theorem.

\section*{Evidence and disclosure}
The reduction proves the theorem independently of computation. The accompanying finite checks are described in the repository. Gia produced the discovery proof; Claude assisted in earlier writing and review, and OpenAI Codex assisted in the present audit and preparation under human direction. Internally reviewed, not peer reviewed.
\bibliographystyle{amsplain}
\bibliography{references}
\end{document}

